\section{Findings}
\subsection{Corpus-Level Trends Before the RQ Analysis}
Before answering the four research questions directly, three corpus-level observations are worth emphasizing.

First, the field is now primarily a systems literature rather than a pure model-scaling literature. Although a few studies still compare standalone small and large models directly, the dominant trend is to embed them inside a larger control strategy. Hybrid edge-cloud pipelines, retrieval-centered systems, expert routers, task decomposers, self-reflective workflows, and fusion modules are now common design choices. In other words, recent work increasingly treats model scale as one variable in a broader systems problem.

Second, hybridization is motivated by different constraints in different domains. In telecom and networking, the main concerns are deployment, context fidelity, and latency budgets. In healthcare and scientific tasks, the dominant concerns are reliability, interpretability, and domain specificity. In industrial monitoring, multimodal hybridization is driven by the need to combine coarse global understanding with fine-grained local detail. In optimisation and education, hybrid methods often serve as a way to recover reasoning quality from small models through feedback, routing, or targeted prompting.

Third, the evidence base is uneven. Accuracy is almost always reported; compute, latency, and cost are not. This matters because a large share of the literature frames hybridization as an efficiency solution, yet only a minority of studies expose sufficiently comparable numbers to test those claims rigorously. The result is a literature that often proves that hybridization \emph{can} work, but less often proves how much it costs in practice under standardized conditions.

\subsection{RQ1: Comparative Performance of Hybrid and Multi-Agent Systems}
Across the corpus, hybrid and multi-agent architectures generally improve the performance ceiling of SLM-centered systems and, in several cases, also outperform larger monolithic baselines. The main reason is not that small models suddenly become universally strong, but that coordination mechanisms route the right subtasks, tokens, or contexts to the right component at the right moment.

The clearest evidence comes from edge-cloud collaboration. Hao et al.~\cite{hao2024hybrid} show that token-level draft-verification can recover LLM-like quality on GSM8K while using only 25.8\% of the full-cloud LLM cost. Their abstract and introduction clarify why this works: many tokens are easy enough for the on-device SLM, and only a small subset of difficult tokens must be verified or regenerated in the cloud. The key performance implication is that quality recovery does not require whole-query escalation. Token-level selectivity is enough.

This pattern also appears in wireless and mobile settings. Oh et al.~\cite{oh2025uncertainty} frame hybrid inference as a distributed speculative-decoding problem and show that uncertainty-aware skipping can preserve up to 97.54\% of LLM accuracy. Their work is especially important for the present review because it ties performance directly to a measurable systems lever: if the SLM can estimate when its token is likely to be accepted, then the system can skip unnecessary communication and cloud inference without losing much quality.

Routing can do more than merely preserve LLM quality. In dialogue state tracking, OrchestraLLM uses a retrieval-based router to select the most suitable expert among SLM and LLM pools~\cite{lee2024orchestrallm}. The website extraction reports 82.46\% turn-level joint goal accuracy on MultiWOZ and 68.09\% on SGD, outperforming both the LLM and SLM baselines. The corresponding PDF makes the reason more explicit: SLMs and LLMs exhibit complementary strengths, with SLMs handling schema and formatting constraints well and LLMs contributing commonsense and broader contextual reasoning. This is an example of a hybrid system exceeding both ends of the scale spectrum because the system exploits complementarity rather than averaging behavior.

Domain-specific retrieval produces similarly strong results. TeleOracle adapts Phi-2 to telecom question answering by combining fine-tuning, semantic chunking, hybrid retrieval, reranking, and context-window extension~\cite{alabbasi2024teleoracle}. In the reviewed corpus, this system reaches 81.20\% accuracy, which is roughly 30\% above the base Phi-2 model and competitive with much larger models. The PDF-level introduction is helpful here because it explains that telecom concepts can conflict with general-language usage. In such a setting, a small model supported by domain retrieval can outperform a larger but poorly grounded general-purpose model.

The reviewed literature also shows that hybridization is not limited to text-only inference. In EnerSafe-Cap, a lightweight video-language model and a multimodal large model are combined through dual-path sampling and prompt-guided fusion for industrial safety monitoring~\cite{cao2025enersafe}. This multimodal setting matters because it demonstrates that hybrid language-model design is also a coordination strategy over different visual abstractions. The lightweight path captures broad continuity, while the larger model focuses on key fine-grained moments. The performance gains over both SwinBERT and Qwen2-VL-2B suggest that hybridization can improve caption quality when different modalities or sampling strategies emphasize genuinely different information.

Several studies in the corpus further show that performance recovery can come from indirect support rather than direct LLM invocation at inference time. Mixed Distillation uses both Chain-of-Thought and Program-of-Thought supervision to transfer reasoning competence into smaller models~\cite{li2024mixed}. In the reviewed data, the resulting 7B models outperform GPT-3.5-Turbo on SVAMP. Likewise, FIPO does not simply ask an LLM to solve the task directly; instead, it uses feedback-guided prompt optimisation to improve small-model problem formulation in mathematical optimisation~\cite{amarasinghe2025fipo}. The result is a dramatic increase in Qwen-2.5 performance from 3\% to 54\% on LPWP. These examples show that hybridization can occur at the training-time or prompt-evolution level, not only at test-time inference.

Despite these positive results, the corpus does not support a simplistic conclusion that SLMs have ``caught up'' to LLMs in general. The strongest comparative results tend to appear in settings where at least one of the following is true: the task is domain-specific, the context can be retrieved and constrained, the output structure is well defined, the system can decompose the problem into sub-parts, or escalation is possible when uncertainty is high. By contrast, open-ended reasoning, general medical question answering, and long-context understanding remain much harder for small models when those supporting conditions are absent.

\subsection{RQ2: Efficiency, Cost, Latency, and Energy Trade-offs}
The corpus strongly suggests that efficiency is one of the main motivations for SLM--LLM collaboration, but it also reveals that the evidence base is much weaker here than in pure performance evaluation. Out of 54 papers, 51 report task-quality measures, while only 17 report latency, 5 report cost, and 4 report energy. This asymmetry is itself a finding: hybrid systems are often justified as efficient, but only a minority of studies provide directly comparable measurements of that efficiency.

When cost is reported, the numbers are often compelling. Hao et al.~\cite{hao2024hybrid} reduce cloud cost to 25.8\%--31.2\% of full LLM inference across three benchmark settings by escalating only a small subset of difficult tokens. Division-of-Thoughts reports 83.57\% lower API cost while maintaining competitive accuracy on reasoning and web-agent tasks~\cite{shao2025dot}. Token-level routing for edge devices, according to the website extraction, routes fewer than 7\% of tokens to the cloud and claims more than 80\% cost reduction relative to full LLM inference. These numbers point to a consistent design pattern: the largest savings come from selective escalation rather than from shrinking the entire pipeline.

Latency evidence is more complex. U-HLM is one of the clearest papers because it reports both local and remote timing and links them to communication overhead~\cite{oh2025uncertainty}. The SLM takes 24.6 ms and the LLM 104.6 ms in the reported setup, and uncertainty-aware skipping improves throughput by 2.54$\times$. Division-of-Thoughts improves average reasoning time by 66.12\% through sub-task parallelism and selective allocation~\cite{shao2025dot}. In contrast, multimodal collaboration may improve accuracy at the cost of much slower per-instance processing: EnerSafe-Cap reports approximately 19--29 seconds per video. These examples show that ``hybrid'' does not imply low latency by default. The latency effect depends on the granularity of cooperation and the burden of the coordination mechanism.

The reviewed PDFs also help explain why latency behaves differently across systems. In hybrid edge-cloud inference, latency overhead comes from token verification and communication. In U-HLM, the dominant bottleneck is vocabulary-distribution upload and repeated model evaluation. In OrchestraLLM, compute savings are reported in FLOPs rather than milliseconds because the router adds negligible cost relative to LLM execution~\cite{lee2024orchestrallm}. In TeleOracle, the website notes that cross-encoder reranking is the latency bottleneck even though the core generator is small. The important systems lesson is that hybrid design moves cost around; it does not eliminate it.

Energy is the least mature evaluation axis in the corpus. Only four papers report energy-related information at all, and only two provide clearly numeric values. One reviewed paper on energy-per-token argues that a small model using Chain-of-Thought can be less energy-efficient than a larger model without such prompting. Another networking paper reports strong resource savings in near-real-time loops. These sparse reports are enough to show that energy cannot be inferred directly from parameter count, but they are not enough to support a high-confidence cross-study conclusion.

Overall, the efficiency evidence supports a cautious but clear synthesis. Hybridization is most beneficial when it suppresses unnecessary large-model calls, narrows the search space, or avoids sending every token, query, or subtask to the cloud. Efficiency gains are weakest or hardest to interpret when the coordination layer adds expensive retrieval, reranking, debate, or multimodal processing that is not itself measured carefully. A mature literature on SLM--LLM systems will therefore need standardized reporting for latency, cloud usage, API cost, and energy, not just accuracy.

\subsection{RQ3: Orchestration and Routing Mechanisms}
The corpus contains a broad set of orchestration mechanisms, but they can be organized into a small number of recurring families. This taxonomy is important because the strongest papers in the review differ less by absolute model size than by how effectively they divide labor.

\paragraph{Token-level collaboration.}
Token-level orchestration appears in speculative or verification-based systems. Hybrid edge-cloud collaborative inference uses the SLM as a drafter and the LLM as a verifier on hard tokens~\cite{hao2024hybrid}. U-HLM adds an uncertainty-aware skipping mechanism so that even this token-level handshake is avoided whenever the SLM's token is likely to be accepted~\cite{oh2025uncertainty}. Token-level routing is powerful because it preserves fine control over quality, but it can suffer from communication and synchronization overhead if too many tokens are escalated.

\paragraph{Query-level or example-level routing.}
OrchestraLLM exemplifies retrieval-based expert routing at the dialogue level~\cite{lee2024orchestrallm}. Similar logic appears in customer-review analysis and several question-answering systems. These designs are attractive when the system can identify ``easy'' and ``hard'' instances in a semantically meaningful way.

\paragraph{Retrieval-centered orchestration.}
RAG is a major control mechanism in the corpus even when the website classifies a paper as hybrid rather than purely RAG. TeleOracle, RAP-RAG, ColBERT retrieval, and several domain QA systems use retrieval not only to expand context but also to constrain the action space of a smaller model. In these systems, retrieval functions as a coordination layer that substitutes for raw model scale by narrowing ambiguity and grounding the generated answer in domain evidence.

\paragraph{Planner-based decomposition and agentic workflows.}
Task decomposition is increasingly important in 2025 papers. Division-of-Thoughts uses a task decomposer, dependency graph, and plug-and-play adapter to allocate sub-tasks between local and remote models~\cite{shao2025dot}. FIPO uses an agentic workflow that generates structured feedback from multiple perspectives before updating the prompt of the small model~\cite{amarasinghe2025fipo}. Flight-deck support, aviation agents, bioinformatics agents, and resilient task-execution frameworks on the website all follow the same idea: if the task can be decomposed into specialists, small models can be retained in the loop even when a large model remains necessary somewhere in the pipeline.

\paragraph{Fusion and multimodal coordination.}
Not all orchestration is text-only. EnerSafe-Cap combines coarse-grained and fine-grained visual semantics through prompt-guided fusion~\cite{cao2025enersafe}. Molecular-property prediction papers also fuse language-model embeddings with graph-based features. In these settings, orchestration is best understood as structured fusion between heterogeneous experts.

These patterns support a broader interpretation of orchestration. It is not merely an implementation detail added after choosing models. In the reviewed literature, orchestration is often the principal source of improvement. The best systems win because they know when to use the SLM, when to ask for help, when to retrieve evidence, and when to decompose the task. This is why the same small models can perform poorly in monolithic form but competitively when embedded in a better control architecture.

\subsection{RQ4: Domain Suitability of SLM-Centered Systems}
The corpus shows that SLM competitiveness is domain-dependent rather than universal. Small models are most successful when at least one of the following conditions holds: the domain vocabulary is narrow, the task is structurally constrained, domain retrieval can ground the generation process, the output is partially symbolic or template-like, or the system can escalate uncertain cases to a larger model.

\paragraph{Telecom and networking.}
Telecom is one of the clearest success stories for SLM-centered systems. TeleOracle and other reviewed telecom QA papers show that small models can perform well when standards documents are retrieved effectively and the generator is trained or prompted for domain fidelity~\cite{alabbasi2024teleoracle}. The reasons are intuitive: the domain is terminology-heavy, the relevant documents are accessible, and faithfulness matters at least as much as broad generative flexibility.

\paragraph{Healthcare and scientific domains.}
Healthcare results are mixed rather than uniformly positive. On the one hand, the website includes strong performance in radiology report labeling, chest X-ray reporting, VTE identification, multilingual medical summarization, and molecular property prediction. These are tasks with domain structure, labeled outputs, or strong contextual priors. On the other hand, the corpus also contains medical question-answering and diagnosis-generation settings where larger models remain more capable or where reasoning prompts drastically inflate latency. The lesson is not that healthcare favors one model scale universally, but that smaller models work best when the output space is constrained and domain knowledge can be tightly anchored.

\paragraph{Reasoning, optimisation, and education.}
Reasoning-heavy tasks produce some of the most interesting hybrid results. Mixed Distillation improves mathematical reasoning enough for 7B models to surpass GPT-3.5-Turbo on SVAMP~\cite{li2024mixed}. FIPO shows that even very weak baseline SLM performance on optimisation formulation can be transformed through feedback-guided prompt evolution~\cite{amarasinghe2025fipo}. Division-of-Thoughts and question-answering agents demonstrate that decomposition and selective allocation can keep small models viable on reasoning workloads that would be implausible in a single-pass local-only setup~\cite{shao2025dot}. These studies suggest that small models are not inherently incapable of reasoning; rather, they are highly sensitive to scaffolding.

\paragraph{Code, software, and agentic tasks.}
The website includes several code-related or software-oriented studies, including code correction, type and call-graph analysis, offline robotic decision-making, task-executing multi-agent systems, and academic feedback generation. These tasks often reward decomposition, specialist routing, or domain grounding more than sheer model size. Still, several reviewed systems note that performance drops as task complexity increases, especially when the agent must synthesize many steps or maintain long interaction histories.

\paragraph{Multimodal industrial monitoring.}
EnerSafe-Cap illustrates a particularly important case because it is not a conventional text benchmark~\cite{cao2025enersafe}. In industrial video monitoring, the hybrid system outperforms both the lightweight and the larger baseline models because the task benefits from two complementary forms of attention: one to continuity and one to salient detail. This suggests that in multimodal domains, the question is often not whether the small model can replace the large one, but whether each can capture a different slice of the evidence.

\paragraph{Failure modes and limitations.}
The domain analysis also reveals several failure modes. Small models struggle most when the task requires broad open-domain knowledge, long-context synthesis without strong retrieval, or flexible reasoning over ambiguous instructions. They also become brittle when prompt structure matters too much or when coordination modules are themselves inaccurate. A central finding of the review is therefore that SLM suitability should be evaluated as a property of \emph{system design plus domain}, not as a stable property of parameter count alone.

\subsection{Cross-RQ Synthesis}
Viewed together, the four research questions suggest that the literature is converging toward a specific systems principle: hybrid language-model design works best when the large model is treated as a selectively invoked capability rather than the default engine for every token, query, or subtask.

The first implication is architectural. Strong results are associated with explicit control layers such as routers, planners, retrievers, or fusion modules. This means that future progress is unlikely to come only from making SLMs slightly larger or LLMs slightly cheaper. It will also depend on better confidence estimation, task decomposition, retrieval selection, and cross-model scheduling.

The second implication is methodological. Because accuracy dominates reporting, many papers overstate practical efficiency. The corpus already contains cases where stronger orchestration clearly reduces cost or latency, but it also contains cases where coordination shifts the bottleneck rather than removing it. A mature evaluation protocol should therefore report resource metrics as first-class outcomes rather than optional extras.

The third implication is application-oriented. The most successful SLM-centered deployments are not generic substitutes for LLMs. They are domain-adapted systems in which the small model is grounded, specialized, or protected by escalation logic. This is why telecom, structured reasoning, radiology labeling, optimisation formulation, and industrial monitoring are so prominent in the positive findings of the corpus.
